%------------------------------------------------------
% Harvard-style one-page resume -- NGUYEN KHAC BAO
% Target: Fullstack AI Integration Intern
% Compile: pdflatex CV_Nguyen_Khac_Bao_Fullstack_AI_Integration_Intern_CloudRun.tex
%------------------------------------------------------

\documentclass[10pt,a4paper]{article}

\usepackage[T1]{fontenc}
\usepackage[utf8]{inputenc}
\usepackage[a4paper,top=0.38in,bottom=0.38in,left=0.52in,right=0.52in]{geometry}
\usepackage{enumitem}
\usepackage{hyperref}
\usepackage{titlesec}
\usepackage{microtype}
\usepackage{xcolor}

\hypersetup{
    colorlinks=true,
    urlcolor=blue,
    linkcolor=black,
    pdfborder={0 0 0}
}

\titleformat{\section}
    {\normalfont\large\bfseries}
    {}
    {0em}
    {}
    [\vspace{1pt}\titlerule]

\titlespacing*{\section}{0pt}{4pt}{2pt}

\setlength{\parindent}{0pt}
\setlength{\parskip}{0pt}
\pagestyle{empty}
\sloppy

\newenvironment{cvitems}{%
    \begin{itemize}[leftmargin=1.2em,itemsep=0pt,topsep=1pt,parsep=0pt]
}{%
    \end{itemize}
}

\newcommand{\entryhead}[2]{%
    \noindent\textbf{#1}\hfill #2\\[-2pt]
}

\newcommand{\entrysub}[1]{%
    \noindent\textit{#1}\\[-1pt]
}

\begin{document}
\fontsize{9.0}{9.9}\selectfont

\begin{center}
    {\Large\textbf{NGUYEN KHAC BAO}}\\[-1pt]
    Fullstack AI Integration Intern\\[2pt]
    Ho Chi Minh City, Vietnam
    \,\textbar\,
    \href{mailto:nguyenkhacbaowork564@gmail.com}{nguyenkhacbaowork564@gmail.com}
    \,\textbar\,
    (+84) 039 586 4429
    \,\textbar\,
    \href{https://github.com/NguyenKhacBao564}{github.com/NguyenKhacBao564}
\end{center}

\section{CAREER OBJECTIVE}
AI integration intern candidate focused on turning AI models into usable products through backend APIs, RAG pipelines, vector databases, Docker, and cloud deployment. Hands-on with FastAPI, Streamlit, Gemini API, Qdrant Cloud, Cloud SQL, SSE-style streaming, and Cloud Run for deployable AI demos.

\section{EDUCATION}
\entryhead{Posts and Telecommunications Institute of Technology (PTIT)}{2022 -- 2027}
\entrysub{Bachelor of Information Technology}
Relevant Coursework: Machine Learning, Deep Learning, Computer Vision.

\section{TECHNICAL SKILLS}
\textbf{AI Integration:} RAG pipelines, Gemini API, embeddings, prompt engineering, OpenAI-compatible APIs, SSE-style streaming, fallback handling\\
\textbf{Backend / API:} Python, FastAPI, REST API, WebSocket, Pydantic\\
\textbf{Databases:} Qdrant Cloud, MySQL / Cloud SQL, SQL, JSON / JSONL, CSV processing\\
\textbf{Cloud / DevOps:} Docker, Google Cloud Run, Cloud Build, Artifact Registry, Secret Manager, Linux\\
\textbf{ML / CV:} PyTorch, YOLO, OpenCV, OCR, ByteTrack\\
\textbf{Tools:} Git / GitHub, Vast.ai (GPU), Hugging Face, Streamlit, Jupyter / Colab, Postman

\section{LANGUAGES}
Vietnamese: Native \textbar\ English: TOEIC 550/990; technical reading proficiency for documentation, model cards, API references, and engineering articles.

\section{PROJECTS}
\entryhead{Vietnamese Legal RAG Chatbot -- Fullstack AI App on Cloud Run}{2026}
\entrysub{Personal Project \textbar\ Python, FastAPI, Streamlit, Gemini API, Qdrant Cloud, Cloud SQL/MySQL, Docker, Google Cloud Run \textbar\ Source: \href{https://github.com/NguyenKhacBao564/legal_RaG_Chatbot_VietNamese-System}{GitHub} \textbar\ Demo: Cloud Run frontend/API health check}
\begin{cvitems}
    \item Integrated Gemini as the default LLM and embedding provider, with Qdrant Cloud retrieval over a Vietnamese legal QA dataset for context-grounded answers.
    \item Built a cloud-ready fullstack RAG chatbot with a Streamlit frontend and FastAPI backend deployed as separate Google Cloud Run services.
    \item Used Cloud SQL/MySQL for chat history and Secret Manager/environment-based configuration for API keys, Qdrant credentials, and database credentials.
    \item Added SSE-style streaming response support, readiness checks, structured latency logging, and fallback handling for Qdrant/LLM failures without breaking the existing synchronous chat endpoint.
\end{cvitems}

\entryhead{Real-Time Face Mask Compliance Detection -- Deployed AI Inference Service}{2026}
\entrysub{Personal Project \textbar\ Python, FastAPI, WebSocket, YOLOv8n, OpenCV, Docker, Google Cloud Run \textbar\ Source: \href{https://github.com/NguyenKhacBao564/Real-Time-Face-Mask-Compliance-Detection-System}{GitHub} \textbar\ Demo: Cloud Run webcam/API demo}
\begin{cvitems}
    \item Built and deployed a real-time AI inference service with FastAPI REST + WebSocket endpoints and a browser webcam frontend running over HTTPS/WSS.
    \item Trained a YOLOv8n 3-class mask detector on PWMFD converted from VOC to YOLO format: 9,205 images, 18,528 labeled instances, and 0.970 mAP@50 validation.
    \item Logged warm endpoint latency around 151 ms for REST and 127 ms for WSS in the demo environment, with lightweight inference/FPS monitoring.
    \item Added violation event logging with per-event metadata (timestamp, confidence, snapshot) and a review API -- mirroring camera-side alerting architecture.
\end{cvitems}

\entryhead{Traffic Violation Detection \& License Plate Recognition -- AI Camera Evidence API}{2026}
\entrysub{Personal Project \textbar\ Python, FastAPI, OpenCV, Ultralytics YOLO, ByteTrack, HyperLPR3, JSON \textbar\ Source: \href{https://github.com/NguyenKhacBao564/Traffic-Violation-Detection-and-License-Plate-Recognition}{GitHub}}
\begin{cvitems}
    \item Built a fixed-camera video analytics pipeline for vehicle detection/tracking, traffic-light ROI estimation, stop-line crossing logic, plate detection/OCR, and JSON evidence bundles.
    \item Built a FastAPI REST API exposing video analysis, event retrieval, and frame/plate evidence endpoints -- enabling programmatic access to pipeline results for downstream integration.
    \item Prepared a compact CCPD2019 plate dataset with annotation parsing, YOLO bbox conversion, train/val/test split, negative samples, OCR crops, and QA reports.
    \item Fine-tuned a YOLO plate detector on 8,598 images, achieving 0.989 precision, 0.998 recall, and 0.994 mAP50 on validation; documented current OCR limitations and local demo constraints.
\end{cvitems}

\section{ACTIVITIES}
\entryhead{AI Research Group -- University Club}{2025 -- Present}
\entrysub{Member}
Discussed Machine Learning, Deep Learning, Computer Vision, RAG/LLM applications, and AI system development.

\end{document}
